\documentclass{article}
\usepackage[utf8]{inputenc}
\usepackage{geometry}
\usepackage{enumitem}
\usepackage{hyperref}
\usepackage{xcolor}
\usepackage{graphicx}
\usepackage{array}
\usepackage{multirow}
\usepackage{amsmath, amssymb}
\usepackage{chemfig}
\usepackage{mhchem}
\usepackage{tikz}
\usepackage{setspace}
\usepackage{soul}
\usepackage{mathtools}
\usepackage{booktabs}
\usepackage{chemformula}
\usepackage{siunitx}
\usetikzlibrary{positioning, calc}
\usepackage{longtable}
\geometry{margin=1in}
\setlength{\parskip}{0.5em}
\usepackage{cancel}


\begin{document}

\begin{center}
    \vspace*{-1cm}
    \begin{tabular}{p{12cm} r}
        & \textbf{Name:} \underline{\hspace{5cm}} \\
    \end{tabular}

    \vspace{0.2cm}
    {\LARGE \textbf{Week 11 Worksheet}}
\end{center}    
\\

\noindent\begin{tabular}{|p{0.9\linewidth}|}
\hline
\textbf{(1)} Why is mass/volume percent often used in medical solutions such as saline?\\[0.2cm]
\begin{minipage}[t]{0.85\linewidth}
\begin{enumerate}[label=\alph*.]
\item Because volume measurements are easier in solids
\item Because both solute and solvent are liquids
\item Because solute is solid and solvent is liquid
\item Because it gives the total molarity directly
\end{enumerate}
\end{minipage}\\[0.2cm]
\hline
\end{tabular}

\noindent\begin{tabular}{|p{0.9\linewidth}|}
\hline
\textbf{(2)} What is the molarity of a solution made by dissolving 49.0~g of H$_2$SO$_4$ in 1.00~L of solution?\\[0.2cm]
\begin{minipage}[t]{0.85\linewidth}
\begin{enumerate}[label=\alph*.]
\item 0.250~M
\item 0.500~M
\item 1.00~M
\item 2.00~M
\end{enumerate}
\end{minipage}\\[0.2cm]
\hline
\end{tabular}

\noindent\begin{tabular}{|p{0.9\linewidth}|}
\hline
\textbf{(3)} What term describes the process by which a solid dissolves until no more can dissolve under given conditions?\\[0.2cm]
\begin{minipage}[t]{0.85\linewidth}
\begin{enumerate}[label=\alph*.]
\item Crystallization
\item Saturation
\item Dilution
\item Precipitation
\end{enumerate}
\end{minipage}\\[0.2cm]
\hline
\end{tabular}

\noindent\begin{tabular}{|p{0.9\linewidth}|}
\hline
\textbf{(4)} A 5\% (m/v) glucose solution contains 5~g of glucose in:\\[0.2cm]
\begin{minipage}[t]{0.85\linewidth}
\begin{enumerate}[label=\alph*.]
\item 5~mL of water
\item 100~mL of water
\item 100~mL of solution
\item 95~mL of solution
\end{enumerate}
\end{minipage}\\[0.2cm]
\hline
\end{tabular}
\newpage

\noindent\begin{tabular}{|p{0.9\linewidth}|}
\hline
\textbf{(5)} If 75.0~mL of 4.00~M KOH are diluted to 300.0~mL, what is the new molarity?\\[0.2cm]
\begin{minipage}[t]{0.85\linewidth}
\begin{enumerate}[label=\alph*.]
\item 0.50~M
\item 1.00~M
\item 2.00~M
\item 4.00~M
\end{enumerate}
\end{minipage}\\[0.2cm]
\hline
\end{tabular}

\noindent\begin{tabular}{|p{0.9\linewidth}|}
\hline
\textbf{(6)} A chemist dissolves 25~g of NaCl in 75~g of water. What is the percent by mass of NaCl?\\[0.2cm]
\begin{minipage}[t]{0.85\linewidth}
\begin{enumerate}[label=\alph*.]
\item 20\%
\item 25\%
\item 30\%
\item 33\%
\end{enumerate}
\end{minipage}\\[0.2cm]
\hline
\end{tabular}

\noindent\begin{tabular}{|p{0.9\linewidth}|}
\hline
\textbf{(7)} How many milliliters of 0.250~M HNO$_3$ are needed to supply 0.0500~mol of HNO$_3$?\\[0.2cm]
\begin{minipage}[t]{0.85\linewidth}
\begin{enumerate}[label=\alph*.]
\item 50.0~mL
\item 100.~mL
\item 150.~mL
\item 200.~mL
\end{enumerate}
\end{minipage}\\[0.2cm]
\hline
\end{tabular}

\noindent\begin{tabular}{|p{0.9\linewidth}|}
\hline
\textbf{(8)} Which of the following would double the (m/v)\% of a solution?\\[0.2cm]
\begin{minipage}[t]{0.85\linewidth}
\begin{enumerate}[label=\alph*.]
\item Double solute, same volume
\item Double volume, same solute
\item Halve solute, double volume
\item Keep solute and volume constant
\end{enumerate}
\end{minipage}\\[0.2cm]
\hline
\end{tabular}
\newpage

\noindent\begin{tabular}{|p{0.9\linewidth}|}
\hline
\textbf{(9)} Why are percent by mass and percent by volume not interchangeable?\\[0.2cm]
\begin{minipage}[t]{0.85\linewidth}
\begin{enumerate}[label=\alph*.]
\item They are measured with different instruments
\item One involves solids while the other involves liquids
\item They describe different types of proportional relationships
\item They are only used for gases
\end{enumerate}
\end{minipage}\\[0.2cm]
\hline
\end{tabular}

\noindent\begin{tabular}{|p{0.9\linewidth}|}
\hline
\textbf{(10)} What does a 15\% (m/v) H\textsubscript{2}O\textsubscript{2} solution contain?\\[0.2cm]
\begin{minipage}[t]{0.85\linewidth}
\begin{enumerate}[label=\alph*.]
\item 15~mL solute in 100~mL solvent
\item 15~g solute in 100~g solution
\item 15~g solute in 100~mL solution
\item 15~mL solute in 100~g solution
\end{enumerate}
\end{minipage}\\[0.2cm]
\hline
\end{tabular}

\noindent\begin{tabular}{|p{0.9\linewidth}|}
\hline
\textbf{(11)} A 25.0~mL sample of 0.100~M Ba(OH)$_2$ neutralizes how many moles of HCl?\\[0.2cm]
\[
\text{Unbalanced chemical equation: } \mathrm{Ba(OH)_2_{(aq)} + HCl_{(aq)} \rightarrow BaCl_2_{(aq)} + H_2O_{(l)}}
\]
\begin{minipage}[t]{0.85\linewidth}
\begin{enumerate}[label=\alph*.]
\item 0.00125~mol
\item 0.00250~mol
\item 0.00500~mol
\item 0.0100~mol
\end{enumerate}
\end{minipage}\\[0.2cm]
\hline
\end{tabular}

\noindent\begin{tabular}{|p{0.9\linewidth}|}
\hline
\textbf{(12)} If additional solute is added to a saturated solution, what happens?\\[0.2cm]
\begin{minipage}[t]{0.85\linewidth}
\begin{enumerate}[label=\alph*.]
\item It dissolves completely
\item It crystallizes out
\item It remains undissolved
\item It lowers the saturation point
\end{enumerate}
\end{minipage}\\[0.2cm]
\hline
\end{tabular}
\newpage

\noindent\begin{tabular}{|p{0.9\linewidth}|}
\hline
\textbf{(13)} Volume/volume percent compares the volume of solute to the volume of:\\[0.2cm]
\begin{minipage}[t]{0.85\linewidth}
\begin{enumerate}[label=\alph*.]
\item Solvent
\item Solution
\item Solute
\item Both solute and solvent
\end{enumerate}
\end{minipage}\\[0.2cm]
\hline
\end{tabular}

\noindent\begin{tabular}{|p{0.9\linewidth}|}
\hline
\textbf{(14)} Which of the following is a common use of (v/v)\%?\\[0.2cm]
\begin{minipage}[t]{0.85\linewidth}
\begin{enumerate}[label=\alph*.]
\item Sugar solutions
\item Alcoholic beverages
\item Saline solutions
\item Solid mixtures
\end{enumerate}
\end{minipage}\\[0.2cm]
\hline
\end{tabular}

\noindent\begin{tabular}{|p{0.9\linewidth}|}
\hline
\textbf{(15)} Which of the following best describes a saturated solution at equilibrium?\\[0.2cm]
\begin{minipage}[t]{0.85\linewidth}
\begin{enumerate}[label=\alph*.]
\item Solute is dissolving faster than it crystallizes
\item Solute and solvent are reacting chemically
\item Dissolving and crystallization occur at equal rates
\item No solute particles are moving
\end{enumerate}
\end{minipage}\\[0.2cm]
\hline
\end{tabular}

\noindent\begin{tabular}{|p{0.9\linewidth}|}
\hline
\textbf{(16)} What volume of 0.500~M NaOH is required to neutralize 25.0~mL of 0.250~M HCl?\\[0.2cm]
\[
\text{Unbalanced chemical equation: } \mathrm{HCl_{(aq)} + NaOH_{(aq)} \rightarrow NaCl_{(aq)} + H_2O_{(l)}}
\]
\begin{minipage}[t]{0.85\linewidth}
\begin{enumerate}[label=\alph*.]
\item 12.5~mL
\item 25.0~mL
\item 50.0~mL
\item 100.0~mL
\end{enumerate}
\end{minipage}\\[0.2cm]
\hline
\end{tabular}
\newpage

\noindent\begin{tabular}{|p{0.9\linewidth}|}
\hline
\textbf{(17)} In mass/volume percent, the solute is measured in \_\_\_\_ and the solution in \_\_\_\_.\\[0.2cm]
\begin{minipage}[t]{0.85\linewidth}
\begin{enumerate}[label=\alph*.]
\item grams; grams
\item grams; liters
\item grams; milliliters
\item milliliters; milliliters
\end{enumerate}
\end{minipage}\\[0.2cm]
\hline
\end{tabular}

\noindent\begin{tabular}{|p{0.9\linewidth}|}
\hline
\textbf{(18)} If 10.0~mL of 12.0~M HCl is diluted to 250.0~mL, what is the final molarity?\\[0.2cm]
\begin{minipage}[t]{0.85\linewidth}
\begin{enumerate}[label=\alph*.]
\item 0.24~M
\item 0.48~M
\item 1.2~M
\item 4.8~M
\end{enumerate}
\end{minipage}\\[0.2cm]
\hline
\end{tabular}

\noindent\begin{tabular}{|p{0.9\linewidth}|}
\hline
\textbf{(19)} A 25\% (v/v) ethanol solution means:\\[0.2cm]
\begin{minipage}[t]{0.85\linewidth}
\begin{enumerate}[label=\alph*.]
\item 25~mL ethanol in 100~mL solution
\item 25~mL ethanol in 100~mL solvent
\item 25~g ethanol in 100~mL solution
\item 25~mL solution in 100~mL ethanol
\end{enumerate}
\end{minipage}\\[0.2cm]
\hline
\end{tabular}

\noindent\begin{tabular}{|p{0.9\linewidth}|}
\hline
\textbf{(20)} A solution that can dissolve more solute at a given temperature is called a(n):\\[0.2cm]
\begin{minipage}[t]{0.85\linewidth}
\begin{enumerate}[label=\alph*.]
\item saturated solution
\item supersaturated solution
\item unsaturated solution
\item concentrated solution
\end{enumerate}
\end{minipage}\\[0.2cm]
\hline
\end{tabular}
\newpage

\noindent\begin{tabular}{|p{0.9\linewidth}|}
\hline
\textbf{(21)} Which of the following is a correct unit for percent by volume concentration?\\[0.2cm]
\begin{minipage}[t]{0.85\linewidth}
\begin{enumerate}[label=\alph*.]
\item g/L
\item mol/L
\item mL solute per 100~mL solution
\item g solute per 100~mL solvent
\end{enumerate}
\end{minipage}\\[0.2cm]
\hline
\end{tabular}

\noindent\begin{tabular}{|p{0.9\linewidth}|}
\hline
\textbf{(22)} A student dissolves 5.85~g of NaCl in enough water to make 500.0~mL of solution. What is the molarity? \\[0.2cm]
\begin{minipage}[t]{0.85\linewidth}
\begin{enumerate}[label=\alph*.]
\item 0.050~M
\item 0.10~M
\item 0.200~M
\item 0.300~M
\end{enumerate}
\end{minipage}\\[0.2cm]
\hline
\end{tabular}

\noindent\begin{tabular}{|p{0.9\linewidth}|}
\hline
\textbf{(23)} What volume of 1.50~M NaOH contains 0.750~mol NaOH?\\[0.2cm]
\begin{minipage}[t]{0.85\linewidth}
\begin{enumerate}[label=\alph*.]
\item 0.125~L
\item 0.250~L
\item 0.500~L
\item 1.00~L
\end{enumerate}
\end{minipage}\\[0.2cm]
\hline
\end{tabular}

\noindent\begin{tabular}{|p{0.9\linewidth}|}
\hline
\textbf{(24)} Which type of percent concentration compares solute mass to the total mass of the solution?\\[0.2cm]
\begin{minipage}[t]{0.85\linewidth}
\begin{enumerate}[label=\alph*.]
\item Percent by volume
\item Percent by mass
\item Percent by molarity
\item Percent by density
\end{enumerate}
\end{minipage}\\[0.2cm]
\hline
\end{tabular}
\newpage

\noindent\begin{tabular}{|p{0.9\linewidth}|}
\hline
\textbf{(25)} If a solution is labeled 70\% (v/v) isopropanol, this means:\\[0.2cm]
\begin{minipage}[t]{0.85\linewidth}
\begin{enumerate}[label=\alph*.]
\item 70~mL of isopropanol in 30~mL of water
\item 70~mL of isopropanol in 100~mL of solution
\item 70~g of isopropanol in 100~g of solution
\item 70~mL of solution in 100~mL of isopropanol
\end{enumerate}
\end{minipage}\\[0.2cm]
\hline
\end{tabular}

\noindent\begin{tabular}{|p{0.9\linewidth}|}
\hline
\textbf{(26)} Molarity, mass percent, volume percent, and mass/volume percent are all units used to express the concentration of a solution.\\[0.2cm]
\begin{minipage}[t]{0.85\linewidth}
\begin{enumerate}[label=\alph*.]
\item True
\item False
\end{enumerate}
\end{minipage}\\[0.2cm]
\hline
\end{tabular}

\noindent\begin{tabular}{|p{0.9\linewidth}|}
\hline
\textbf{(27)} What happens to the molarity of a solution when it is diluted with more solvent?\\[0.2cm]
\begin{minipage}[t]{0.85\linewidth}
\begin{enumerate}[label=\alph*.]
\item It increases
\item It decreases
\item It stays the same
\item It doubles
\end{enumerate}
\end{minipage}\\[0.2cm]
\hline
\end{tabular}

\noindent\begin{tabular}{|p{0.9\linewidth}|}
\hline
\textbf{(28)} When preparing a solution of known concentration, why must all solute be completely dissolved before dilution to the mark?\\[0.2cm]
\begin{minipage}[t]{0.85\linewidth}
\begin{enumerate}[label=\alph*.]
\item To ensure accurate final concentration
\item To allow faster mixing
\item To prevent supersaturation
\item To remove impurities
\end{enumerate}
\end{minipage}\\[0.2cm]
\hline
\end{tabular}
\newpage

\noindent\begin{tabular}{|p{0.9\linewidth}|}
\hline
\textbf{(29)} Volume percent (\% v/v) expresses the volume of solute in milliliters per 100 milliliters of solution.\\[0.2cm]
\begin{minipage}[t]{0.85\linewidth}
\begin{enumerate}[label=\alph*.]
\item True
\item False
\end{enumerate}
\end{minipage}\\[0.2cm]
\hline
\end{tabular}

\noindent\begin{tabular}{|p{0.9\linewidth}|}
\hline
\textbf{(30)} A 10\% (m/m) NaCl solution contains 10 grams of NaCl in 100 grams of total solution (solute + solvent).\\[0.2cm]
\begin{minipage}[t]{0.85\linewidth}
\begin{enumerate}[label=\alph*.]
\item True
\item False
\end{enumerate}
\end{minipage}\\[0.2cm]
\hline
\end{tabular}

\noindent\begin{tabular}{|p{0.9\linewidth}|}
\hline
\textbf{(31)} Mass/volume percent expresses the mass of solute in grams per 100 milliliters of solution.\\[0.2cm]
\begin{minipage}[t]{0.85\linewidth}
\begin{enumerate}[label=\alph*.]
\item True
\item False
\end{enumerate}
\end{minipage}\\[0.2cm]
\hline
\end{tabular}

\noindent\begin{tabular}{|p{0.9\linewidth}|}
\hline
\textbf{(32)} Percent by mass and percent by volume can always be used interchangeably for any solution.\\[0.2cm]
\begin{minipage}[t]{0.85\linewidth}
\begin{enumerate}[label=\alph*.]
\item True
\item False
\end{enumerate}
\end{minipage}\\[0.2cm]
\hline
\end{tabular}
\newpage

\vspace{0.2cm}
\noindent\begin{tabular}{|p{0.9\linewidth}|}
\hline
\textbf{(33)} How many moles of NaCl are present in 250.0~mL of 0.200~M NaCl solution?\\[0.2cm]
\begin{minipage}[t]{0.85\linewidth}
\begin{enumerate}[label=\alph*.]
\item 0.00500~mol
\item 0.0500~mol
\item 0.500~mol
\item 5.00~mol
\end{enumerate}
\end{minipage}\\[0.2cm]
\hline
\end{tabular}

\noindent\begin{tabular}{|p{0.9\linewidth}|}
\hline
\textbf{(34)} What is the molarity of a solution containing 0.250~mol solute in 500.0~mL of solution?\\[0.2cm]
\begin{minipage}[t]{0.85\linewidth}
\begin{enumerate}[label=\alph*.]
\item 0.125~M
\item 0.25~M
\item 0.50~M
\item 1.00~M
\end{enumerate}
\end{minipage}\\[0.2cm]
\hline
\end{tabular}

\noindent\begin{tabular}{|p{0.9\linewidth}|}
\hline
\textbf{(35)} A student heats a saturated solution and all the solid dissolves. When cooled, crystals begin to form. What does this indicate?\\[0.2cm]
\begin{minipage}[t]{0.85\linewidth}
\begin{enumerate}[label=\alph*.]
\item The solubility decreases with temperature
\item The solubility increases with temperature
\item The solvent evaporated
\item The solution was unsaturated
\end{enumerate}
\end{minipage}\\[0.2cm]
\hline
\end{tabular}

\noindent\begin{tabular}{|p{0.9\linewidth}|}
\hline
\textbf{(36)} What is the (v/v)\% of a solution with 40~mL ethanol and total volume 200~mL?\\[0.2cm]
\begin{minipage}[t]{0.85\linewidth}
\begin{enumerate}[label=\alph*.]
\item 10\%
\item 15\%
\item 20\%
\item 25\%
\end{enumerate}
\end{minipage}\\[0.2cm]
\hline
\end{tabular}
\newpage

\noindent\begin{tabular}{|p{0.9\linewidth}|}
\hline
\textbf{(37)} Which of the following describes a dilute saturated solution?\\[0.2cm]
\begin{minipage}[t]{0.85\linewidth}
\begin{enumerate}[label=\alph*.]
\item Contains little solute but cannot dissolve more
\item Contains a large amount of solute and can dissolve more
\item Is very concentrated and still dissolving
\item Is both supersaturated and concentrated
\end{enumerate}
\end{minipage}\\[0.2cm]
\hline
\end{tabular}

\noindent\begin{tabular}{|p{0.9\linewidth}|}
\hline
\textbf{(38)} A 250~mL solution contains 12.5~g of solute. What is its (m/v)\%?\\[0.2cm]
\begin{minipage}[t]{0.85\linewidth}
\begin{enumerate}[label=\alph*.]
\item 2.0\%
\item 5.0\%
\item 6.25\%
\item 10.0\%
\end{enumerate}
\end{minipage}\\[0.2cm]
\hline
\end{tabular}

\noindent\begin{tabular}{|p{0.9\linewidth}|}
\hline
\textbf{(39)} Which type of solution is most unstable?\\[0.2cm]
\begin{minipage}[t]{0.85\linewidth}
\begin{enumerate}[label=\alph*.]
\item Saturated
\item Unsaturated
\item Supersaturated
\item Dilute
\end{enumerate}
\end{minipage}\\[0.2cm]
\hline
\end{tabular}

\noindent\begin{tabular}{|p{0.9\linewidth}|}
\hline
\textbf{(40)} A 12\% (v/v) acetic acid solution means that the solution contains:\\[0.2cm]
\begin{minipage}[t]{0.85\linewidth}
\begin{enumerate}[label=\alph*.]
\item 12~mL acetic acid per 100~mL solution
\item 12~g acetic acid per 100~mL solution
\item 12~mL acetic acid per 100~mL solvent
\item 12~g acetic acid per 100~g solvent
\end{enumerate}
\end{minipage}\\[0.2cm]
\hline
\end{tabular}
\newpage

\noindent\begin{tabular}{|p{0.9\linewidth}|}
\hline
\textbf{(41)} A 100.0~mL sample of 1.20~M Na$_2$CO$_3$ is diluted to 400.0~mL. What is the new concentration?\\[0.2cm]
\begin{minipage}[t]{0.85\linewidth}
\begin{enumerate}[label=\alph*.]
\item 0.30~M
\item 0.60~M
\item 1.20~M
\item 2.40~M
\end{enumerate}
\end{minipage}\\[0.2cm]
\hline
\end{tabular}

\noindent\begin{tabular}{|p{0.9\linewidth}|}
\hline
\textbf{(42)} Which of the following will most likely increase the solubility of a solid in water?\\[0.2cm]
\begin{minipage}[t]{0.85\linewidth}
\begin{enumerate}[label=\alph*.]
\item Decreasing temperature
\item Increasing temperature
\item Decreasing pressure
\item Adding more solute
\end{enumerate}
\end{minipage}\\[0.2cm]
\hline
\end{tabular}

\noindent\begin{tabular}{|p{0.9\linewidth}|}
\hline
\textbf{(43)} A 10\% (m/m) salt solution means there are 10~g of salt in:\\[0.2cm]
\begin{minipage}[t]{0.85\linewidth}
\begin{enumerate}[label=\alph*.]
\item 10~g of solvent
\item 100~g of solvent
\item 100~g of solution
\item 10~mL of solution
\end{enumerate}
\end{minipage}\\[0.2cm]
\hline
\end{tabular}

\noindent\begin{tabular}{|p{0.9\linewidth}|}
\hline
\textbf{(44)} What is the correct formula for percent by volume?\\[0.2cm]
\begin{minipage}[t]{0.85\linewidth}
\begin{enumerate}[label=\alph*.]
\item $\dfrac{\text{volume of solute}}{\text{volume of solution}} \times 100$
\item $\dfrac{\text{mass of solute}}{\text{volume of solvent}} \times 100$
\item $\dfrac{\text{volume of solute}}{\text{mass of solvent}} \times 100$
\item $\dfrac{\text{mass of solute}}{\text{mass of solution}} \times 100$
\end{enumerate}
\end{minipage}\\[0.2cm]
\hline
\end{tabular}
\newpage

\noindent\begin{tabular}{|p{0.9\linewidth}|}
\hline
\textbf{(45)} When a solid is added to water and some dissolves while some remains at the bottom, the solution is:\\[0.2cm]
\begin{minipage}[t]{0.85\linewidth}
\begin{enumerate}[label=\alph*.]
\item Unsaturated
\item Supersaturated
\item Saturated
\item Dilute
\end{enumerate}
\end{minipage}\\[0.2cm]
\hline
\end{tabular}

\noindent\begin{tabular}{|p{0.9\linewidth}|}
\hline
\textbf{(46)} When no more solute can dissolve in a solvent at a given temperature, the solution is said to be:\\[0.2cm]
\begin{minipage}[t]{0.85\linewidth}
\begin{enumerate}[label=\alph*.]
\item unsaturated
\item saturated
\item supersaturated
\item dilute
\end{enumerate}
\end{minipage}\\[0.2cm]
\hline
\end{tabular}

\noindent\begin{tabular}{|p{0.9\linewidth}|}
\hline
\textbf{(47)} A chemist has 50.0~mL of 2.0~M HCl solution. She adds water to dilute it to a final volume of 200.0~mL. What is the molarity of the diluted solution?\\[0.2cm]
\begin{minipage}[t]{0.85\linewidth}
\begin{enumerate}[label=\alph*.]
\item 0.25~M
\item 0.50~M
\item 1.0~M
\item 2.0~M
\end{enumerate}
\end{minipage}\\[0.2cm]
\hline
\end{tabular}

\noindent\begin{tabular}{|p{0.9\linewidth}|}
\hline
\textbf{(48)} The formula for mass/volume percent concentration is:\\[0.2cm]
\begin{minipage}[t]{0.85\linewidth}
\begin{enumerate}[label=\alph*.]
\item $\dfrac{\text{mass of solute}}{\text{mass of solution}} \times 100$
\item $\dfrac{\text{mass of solute}}{\text{volume of solvent}} \times 100$
\item $\dfrac{\text{mass of solute}}{\text{volume of solution}} \times 100$
\item $\dfrac{\text{volume of solute}}{\text{mass of solution}} \times 100$
\end{enumerate}
\end{minipage}\\[0.2cm]
\hline
\end{tabular}
\newpage

\noindent\begin{tabular}{|p{0.9\linewidth}|}
\hline
\textbf{(49)} What is the (m/v)\% of a solution with 8.0~g of KCl in 160~mL of solution?\\[0.2cm]
\begin{minipage}[t]{0.85\linewidth}
\begin{enumerate}[label=\alph*.]
\item 2.0\%
\item 4.0\%
\item 5.0\%
\item 8.0\%
\end{enumerate}
\end{minipage}\\[0.2cm]
\hline
\end{tabular}

\noindent\begin{tabular}{|p{0.9\linewidth}|}
\hline
\textbf{(50)} How many liters of 3.00~M H$_2$SO$_4$ contain 6.00~mol of H$_2$SO$_4$?\\[0.2cm]
\begin{minipage}[t]{0.85\linewidth}
\begin{enumerate}[label=\alph*.]
\item 0.500~L
\item 1.00~L
\item 2.00~L
\item 3.00~L
\end{enumerate}
\end{minipage}\\[0.2cm]
\hline
\end{tabular}

\noindent\begin{tabular}{|p{0.9\linewidth}|}
\hline
\textbf{(51)} Percent concentration expresses the amount of solute in a given amount of solution as a percentage of:\\[0.2cm]
\begin{minipage}[t]{0.85\linewidth}
\begin{enumerate}[label=\alph*.]
\item Only the solute itself (ignores the solvent)
\item Only the solvent (ignores the solute)
\item The total solution (solute + solvent combined; this is what is used in the denominator to calculate percent)
\item The number of moles of solute
\end{enumerate}
\end{minipage}\\[0.2cm]
\hline
\end{tabular}

\noindent\begin{tabular}{|p{0.9\linewidth}|}
\hline
\textbf{(52)} How many grams of KNO$_3$ are required to prepare 250.0~mL of 0.200~M KNO$_3$ solution?\\[0.2cm]
\begin{minipage}[t]{0.85\linewidth}
\begin{enumerate}[label=\alph*.]
\item 2.53~g
\item 5.06~g
\item 10.1~g
\item 20.2~g
\end{enumerate}
\end{minipage}\\[0.2cm]
\hline
\end{tabular}
\newpage

\noindent\begin{tabular}{|p{0.9\linewidth}|}
\hline
\textbf{(53)} How many moles of solute are in 150.0~mL of 0.400~M glucose solution?\\[0.2cm]
\begin{minipage}[t]{0.85\linewidth}
\begin{enumerate}[label=\alph*.]
\item 0.0060~mol
\item 0.0150~mol
\item 0.0600~mol
\item 0.600~mol
\end{enumerate}
\end{minipage}\\[0.2cm]
\hline
\end{tabular}

\noindent\begin{tabular}{|p{0.9\linewidth}|}
\hline
\textbf{(54)} Stirring a solution generally increases the rate of dissolving because it:\\[0.2cm]
\begin{minipage}[t]{0.85\linewidth}
\begin{enumerate}[label=\alph*.]
\item Lowers the solubility
\item Increases surface contact between solute and solvent
\item Changes the solvent's polarity
\item Lowers the temperature
\end{enumerate}
\end{minipage}\\[0.2cm]
\hline
\end{tabular}

\noindent\begin{tabular}{|p{0.9\linewidth}|}
\hline
\textbf{(55)} Why does heating usually increase the solubility of most solid solutes in water?\\[0.2cm]
\begin{minipage}[t]{0.85\linewidth}
\begin{enumerate}[label=\alph*.]
\item It decreases the kinetic energy of particles
\item It allows solvent molecules to move faster and interact more with solute
\item It changes the solvent's polarity
\item It reduces the solute's mass
\end{enumerate}
\end{minipage}\\[0.2cm]
\hline
\end{tabular}

\noindent\begin{tabular}{|p{0.9\linewidth}|}
\hline
\textbf{(56)} Which of the following represents a 10\% (m/v) NaCl solution?\\[0.2cm]
\begin{minipage}[t]{0.85\linewidth}
\begin{enumerate}[label=\alph*.]
\item 10~g NaCl in 100~mL solvent
\item 10~g NaCl in 100~mL solution
\item 10~mL NaCl in 100~mL solution
\item 10~g NaCl in 90~mL water
\end{enumerate}
\end{minipage}\\[0.2cm]
\hline
\end{tabular}
\newpage

\noindent\begin{tabular}{|p{0.9\linewidth}|}
\hline
\textbf{(57)} A solution that contains more dissolved solute than it normally can at that temperature is called:\\[0.2cm]
\begin{minipage}[t]{0.85\linewidth}
\begin{enumerate}[label=\alph*.]
\item saturated
\item unsaturated
\item supersaturated
\item insoluble
\end{enumerate}
\end{minipage}\\[0.2cm]
\hline
\end{tabular}

\noindent\begin{tabular}{|p{0.9\linewidth}|}
\hline
\textbf{(58)} Which of the following statements correctly describes a 5\% (m/m) sugar solution?\\[0.2cm]
\begin{minipage}[t]{0.85\linewidth}
\begin{enumerate}[label=\alph*.]
\item 5~g sugar in 95~g water
\item 5~g sugar in 100~g water
\item 5~g sugar in 100~g solution
\item 5~mL sugar in 100~mL water
\end{enumerate}
\end{minipage}\\[0.2cm]
\hline
\end{tabular}

\noindent\begin{tabular}{|p{0.9\linewidth}|}
\hline
\textbf{(59)} Mass/volume percent is commonly used when the solute is a solid and the solvent is a liquid.\\[0.2cm]
\begin{minipage}[t]{0.85\linewidth}
\begin{enumerate}[label=\alph*.]
\item True
\item False
\end{enumerate}
\end{minipage}\\[0.2cm]
\hline
\end{tabular}

\noindent\begin{tabular}{|p{0.9\linewidth}|}
\hline
\textbf{(60)} In a 20\% (v/v) ethanol solution, 20\% refers to:\\[0.2cm]
\begin{minipage}[t]{0.85\linewidth}
\begin{enumerate}[label=\alph*.]
\item The ratio of solute mass to total mass
\item The ratio of solute volume to total volume
\item The ratio of solvent mass to total mass
\item The ratio of moles of solute to solvent
\end{enumerate}
\end{minipage}\\[0.2cm]
\hline
\end{tabular}

\end{document}
